\label{chap:data}

\section{Sorgenti dei dati}

\subsection{Puzzle Lichess}

La sorgente principale è il database pubblico dei puzzle Lichess, distribuito come
CSV compresso con Zstandard. Ogni record contiene, tra gli altri campi,
l'identificativo del puzzle, una posizione FEN, una sequenza di mosse UCI, il rating
e i temi assegnati. Lichess documenta inoltre una semantica importante: la FEN
precede la mossa dell'avversario che avvia il puzzle, mentre la soluzione del
giocatore comincia dalla seconda mossa della sequenza
\cite{lichessOpenDatabase}.

GraphMate seleziona i temi \texttt{mateIn1}, \texttt{mateIn2},
\texttt{mateIn3}, \texttt{mateIn4} e \texttt{mateIn5}. La pipeline ufficiale
registra 1.937.001 righe dopo il filtraggio tematico e 1.931.434 dopo la pulizia.
Questi valori provengono dai file e dal log della pipeline finale; sostituiscono i
conteggi di una precedente esecuzione esplorativa, non usata per gli split
congelati.

\subsection{Partite ausiliarie e benchmark esterno}

Il progetto include anche un campione di partite Lichess in formato PGN, usato
nella fase esplorativa per studiare i clock e i metadati di gioco. Questi dati non
sono allineati uno-a-uno ai puzzle e non forniscono i tempi usati dal Modello B:
gli esperimenti finali impiegano timing sintetico associato ai puzzle, come
descritto nella Sezione~\ref{sec:synthetic-timing}.

La valutazione esterna usa invece \texttt{yacpdb\_classic\_v1}, un insieme
congelato di 200 composizioni di matto diretto provenienti da YACPDB, con 20
problemi per profondità da \(\#1\) a \(\#10\). Il benchmark è escluso da
addestramento, selezione degli iperparametri e scelta dei checkpoint. La sua
costruzione, la verifica indipendente con Popeye e la politica delle chiavi
accettate sono documentate nell'Appendice~\ref{app:classic-benchmark}.

\section{Pulizia e costruzione del task supervisionato}

La pulizia valida le FEN, richiede rating nell'intervallo 400--3500, profondità
supportata e sequenza non vuota, quindi rimuove i duplicati individuati dalla coppia
FEN--sequenza di mosse. Per costruire l'esempio supervisionato, la pipeline:

\begin{enumerate}
    \item conserva la FEN grezza come \texttt{OriginalFEN};
    \item applica \texttt{Moves[0]} come mossa di preparazione legale;
    \item salva la posizione risultante come FEN del solutore;
    \item usa \texttt{Moves[1]} come \texttt{TargetMove}.
\end{enumerate}

Sono scartati i record con meno di due mosse, FEN non valida, mossa di preparazione
malformata o illegale e target malformato o illegale nella posizione del solutore.
Il compito consiste dunque nel predire la prima mossa risolutiva a partire dalla
posizione effettivamente presentata al giocatore, non nel generare l'intera linea
di matto.

\section{Bilanciamento e suddivisione}

Per limitare lo sbilanciamento della sorgente, vengono campionati al massimo 20.000
esempi per profondità. Le classi mate-in-1--mate-in-4 raggiungono tale soglia,
mentre per mate-in-5 sono disponibili 6.198 esempi. Il dataset bilanciato contiene
quindi 86.198 record.

La suddivisione 80/10/10 usa random state 42 ed è stratificata per profondità. La
Tabella~\ref{tab:main-dataset-splits} distingue le righe CSV dai grafi
effettivamente generati.

\begin{table}[htbp]
    \centering
    \small
    \caption{Split finali e campioni utilizzabili dai modelli a grafo.}
    \label{tab:main-dataset-splits}
    \begin{tabular}{lrrr}
        \toprule
        Split & Righe CSV & Grafi PyG & Target OOV esclusi \\
        \midrule
        Addestramento & 68.958 & 68.958 & 0 \\
        Validazione   & 8.620  & 8.612  & 8 \\
        Test          & 8.620  & 8.610  & 10 \\
        \bottomrule
    \end{tabular}
\end{table}

I primi quattro livelli rappresentano ciascuno circa il 23,20\% di ogni split;
mate-in-5 rappresenta circa il 7,19\%. I rating medi sono rispettivamente
1.338,84, 1.337,61 e 1.338,05 per addestramento, validazione e test. La vicinanza
di queste statistiche è coerente con la suddivisione stratificata, ma non elimina
altre possibili differenze tra esempi.

\section{Codifica della mossa e generazione dei grafi}

Il target resta una stringa UCI fino alla codifica numerica. Il vocabolario è
costruito esclusivamente sull'insieme di addestramento e contiene 1.786 mosse
distinte. Questa scelta evita di usare validazione o test per definire le classi,
ma produce 8 target fuori vocabolario in validazione e 10 nel test. Tali record
rimangono nei CSV e vengono esclusi dalla generazione PyG; per questo le
valutazioni congelate usano 8.610 grafi di test.

Ogni record utilizzabile viene trasformato nel grafo descritto nel
Capitolo~\ref{chap:graph-representation}. I grafi sono salvati in shard da 1.000
elementi e accompagnati da un manifesto che registra schema, conteggi, file e
anomalie. I conteggi dettagliati e il funzionamento del loader sharded sono
riportati nell'Appendice~\ref{app:data-pipeline}.

\section{Variante con timing sintetico}
\label{sec:synthetic-timing}

Il Modello B usa gli stessi esempi e la stessa appartenenza agli split, arricchiti
con il tempo sintetico della mossa precedente, il tempo sintetico della decisione
corrente e un indicatore che ne dichiara la natura sintetica. La generazione è
deterministica rispetto a seed 42, identificativo del puzzle e stream casuale; la
distribuzione è condizionata dal rating e i secondi sono limitati all'intervallo
\([0{,}5,120]\).

La normalizzazione applica uno z-score a \(\log(1+t)\), con media e deviazione
standard stimate sul solo insieme di addestramento. Nessun esempio finale contiene
un tempo umano reale. Questa distinzione è essenziale: l'esperimento misura
l'effetto della specifica codifica sintetica implementata, non l'utilità generale
del tempo di riflessione. Equazioni, parametri e campi persistiti sono riportati
nell'Appendice~\ref{app:data-pipeline}.
